\documentclass[spanish,11pt,a4paper]{article}
\usepackage[margin=1in]{geometry}
\usepackage{fontspec}
\setmainfont{lmroman10-regular.otf}[BoldFont=lmroman10-bold.otf,ItalicFont=lmroman10-italic.otf,BoldItalicFont=lmroman10-bolditalic.otf]
\setsansfont{lmsans10-regular.otf}[BoldFont=lmsans10-bold.otf]
\setmonofont{lmmono10-regular.otf}[Scale=MatchLowercase]
\newfontfamily\LeanSymbols{seguisym.ttf}[Path=typeset_runtime/]
\newlength{\LeanCharWidth}
\usepackage{amsmath,amssymb}
\usepackage[spanish,es-noshorthands,es-tabla]{babel}
\usepackage{microtype,parskip}
\usepackage{graphicx,float,caption}
\usepackage{longtable,booktabs,array,calc}
\usepackage{xcolor}
\usepackage{xurl}
\usepackage{hyperref,bookmark}
\usepackage{etoolbox,needspace}
\setcounter{secnumdepth}{-\maxdimen}
\setlength{\emergencystretch}{3em}
\providecommand{\tightlist}{\setlength{\itemsep}{0pt}\setlength{\parskip}{0pt}}
\newcounter{none}
\AtBeginEnvironment{longtable}{\Needspace{8\baselineskip}\small}
\pretocmd{\section}{\Needspace{5\baselineskip}}{}{}
\pretocmd{\subsection}{\Needspace{4\baselineskip}}{}{}
\pretocmd{\subsubsection}{\Needspace{4\baselineskip}}{}{}
\captionsetup{font=small,labelfont=bf}
\urlstyle{same}
\providecommand{\pandocbounded}[1]{#1}
\hypersetup{pdftitle={Paper I: formalización y reproducibilidad},pdfauthor={Juan Pablo Traverso Gianini},pdflang={es-ES},colorlinks=true,linkcolor=blue,urlcolor=blue,citecolor=blue}
\title{Paper I: formalización y reproducibilidad}
\author{Juan Pablo Traverso Gianini\\\small Investigador independiente, Santiago, Chile\\\small \href{mailto:jtraverso@gmail.com}{jtraverso@gmail.com}\\\small ORCID: 0009-0003-6068-4096}
\date{}
\begin{document}
\raggedbottom
\maketitle
\textbf{Paper I · Formal companion v1.4}\\
19 de octubre de 2026

Este documento acompaña al artículo matemático. Conserva las
declaraciones formales, las precisiones de alcance y la evidencia
computacional que antes figuraban en el artículo. El congelado de
fuentes y su verificación registrada no cambian. No se ejecutó Lean/Lake
para esta edición. Los números históricos de secciones y tablas que
siguen remiten al artículo anterior; se conservan como localizadores y
no se renumeran como resultados nuevos de verificación. Las referencias
originales {[}1{]}, {[}2{]}, etc. figuran al final.

\href{https://github.com/jtraverso/erdos-81-chordal-clique-partitions/tree/eb335845bfa23375d5040ea31fc4f4b366346287/preprints/PAPER_I/05_formalization}{Frozen
sources and historical evidence / Fuentes congeladas y evidencias
históricas} ·
\href{https://github.com/jtraverso/erdos-81-chordal-clique-partitions/tree/eb335845bfa23375d5040ea31fc4f4b366346287/preprints/PAPER_I/01_manuscript}{Previous
article / Artículo anterior}

\subsubsection{Teorema 1.1}\label{teorema-1.1}

\textbf{Estado de la formalización.} El freeze Paper I Lean v1.2 usa
Lean 4 / Mathlib v4.28.0 {[}6,7{]} y registra el Teorema 1.1 como
\nolinkurl{paperI_main_sharp}. La declaración \nolinkurl{paperI_main}
proporciona una forma compañera más débil con término aditivo \(n\). En
estos nombres registrados, \nolinkurl{sharp} se refiere al coeficiente
cuadrático \(1/6\), no a la optimalidad del término aditivo. El gate de
axiomas cubre el teorema principal, \nolinkurl{PaperI.assembly_sharp},
\nolinkurl{PaperI.Split.residual_duality}, la dualidad finita
empaquetamiento--cobertura de la Proposición B.1 y las declaraciones de
interfaz enumeradas en la Sección 10. Su huella registrada es
\nolinkurl{propext}, \nolinkurl{Classical.choice} y \nolinkurl{Quot.sound}, sin
axiomas matemáticos propios del proyecto. La Sección 10 presenta el
registro del freeze y del archivo.

\subsection{10. Reproducibilidad y registro de la
demostración}\label{reproducibilidad-y-registro-de-la-demostraciuxf3n}

La demostración es analítica. Durante el desarrollo se usaron
experimentos computacionales para buscar contraejemplos y detectar
errores de transcripción, pero ningún cálculo es una premisa lógica del
teorema. Por esa razón, el material de regresión computacional se
mantiene fuera del artículo, en lugar de incluirse como un segundo
apéndice.

El ledger de cálculos adjunto registra:

\begin{itemize}
\tightlist
\item
  cada identidad exacta;
\item
  cada dirección de desigualdad;
\item
  la contabilidad completa de pérdidas;
\item
  la rama \(q=0\);
\item
  los regímenes de órbitas;
\item
  la conversión final a \(n\).
\end{itemize}

El autor reporta una auditoría adversarial interna asistida por IA del
enunciado del teorema, la cota cuantitativa y el ensamblaje de la
demostración. Este es un registro de validación, no una afirmación de
revisión externa por pares.

\subsubsection{Estado de formalización en
Lean}\label{estado-de-formalizaciuxf3n-en-lean}

El snapshot en \nolinkurl{PAPER_I/05_formalization/lean_v1.2_freeze}
contiene las fuentes Lean y un build registrado exitoso de 8.034 jobs
con Lean 4 / Mathlib v4.28.0 {[}6,7{]}. Contiene
\nolinkurl{paperI_main_sharp}, cuyo enunciado es
\text{\(\Phi\le n^2/6+n/2\)}, junto con la declaración compañera
más débil \nolinkurl{paperI_main} y los resultados de interfaz enumerados
más abajo. Su gate de axiomas también consulta
\nolinkurl{PaperI.assembly_sharp} y
\nolinkurl{PaperI.Split.residual_duality}. El Apéndice C registra la
declaración principal. Posteriormente, el archivo exacto se reprodujo en
un clean room externo con el mismo resultado de 8.034 jobs y la misma
huella axiomática por teorema.

El informe de axiomas congelado registra únicamente los axiomas
fundacionales estándar de Lean, \nolinkurl{propext},
\nolinkurl{Classical.choice} y \nolinkurl{Quot.sound}, para el teorema
principal y las declaraciones enumeradas; no registra ningún axioma
matemático propio del proyecto. Se trata de fundamentos lógicos usados
en todo Mathlib, no de supuestos matemáticos adicionales del artículo.

El artefacto de release también registra la dualidad finito-dimensional
empaquetamiento-cobertura usada en Proposición B.1. En Lean aparece como
\nolinkurl{residual_duality}, una aplicación del teorema general
\nolinkurl{covering_packing_duality}. El artefacto registra este último
teorema como derivado de un lema de Farkas finito y de la
infraestructura de conos cerrados de Mathlib, en lugar de postularlo
como un axioma externo. Schrijver {[}4, Capítulo 7{]} sigue siendo la
referencia clásica para la demostración escrita.

El archivo de fuentes congelado es
\nolinkurl{PAPER_I_lean_v1.2_freeze.zip}, con SHA-256
\nolinkurl{0181506408644fc1f8872d711de5a98a500f4052aa295bcd6f8c82776694fd3a}.
Se almacena bajo \nolinkurl{PAPER_I/05_formalization/lean_v1.2_freeze}
y excluye los caches del proyecto y los objetos compilados. El target
del build congelado incluye varios módulos \nolinkurl{Contrib} destinados a
un posible upstream; esa vía de contribución es un subproducto y no
forma parte de los supuestos matemáticos de confianza del artículo. El
informe de axiomas anterior controla la afirmación sobre dependencias
lógicas. La reproducción externa en clean room corresponde exactamente a
estos bytes del archivo; la revisión del manuscrito no modifica el
artefacto formal.

\begin{center}\rule{0.5\linewidth}{0.5pt}\end{center}

\subsection{Apéndice C. Enunciado
formalizado}\label{apuxe9ndice-c.-enunciado-formalizado}

El artefacto formal adjunto registra la siguiente declaración Lean 4
para el Teorema 1.1:

\begin{quote}
{\LeanSymbols\fontsize{8}{10}\selectfont
\noindent \nolinkurl{theorem}\hspace{0.45em}\nolinkurl{paperI_main_sharp}\hspace{0.45em}\nolinkurl{(G}\hspace{0.45em}\nolinkurl{:}\hspace{0.45em}\nolinkurl{Split)}\hspace{0.45em}\nolinkurl{:}\par
\noindent \hspace{1.8em}\nolinkurl{G.Phi}\hspace{0.45em}\nolinkurl{≤}\hspace{0.45em}\nolinkurl{(G.n}\hspace{0.45em}\nolinkurl{:}\hspace{0.45em}\nolinkurl{ℝ)}\hspace{0.45em}\nolinkurl{^}\hspace{0.45em}\nolinkurl{2}\hspace{0.45em}\nolinkurl{/}\hspace{0.45em}\nolinkurl{6}\hspace{0.45em}\nolinkurl{+}\hspace{0.45em}\nolinkurl{(G.n}\hspace{0.45em}\nolinkurl{:}\hspace{0.45em}\nolinkurl{ℝ)}\hspace{0.45em}\nolinkurl{/}\hspace{0.45em}\nolinkurl{2}\par
\noindent \par
\noindent \nolinkurl{--}\hspace{0.45em}\nolinkurl{forma}\hspace{0.45em}\nolinkurl{compañera}\hspace{0.45em}\nolinkurl{más}\hspace{0.45em}\nolinkurl{débil}\hspace{0.45em}\nolinkurl{(término}\hspace{0.45em}\nolinkurl{aditivo}\hspace{0.45em}\nolinkurl{mayor):}\par
\noindent \nolinkurl{theorem}\hspace{0.45em}\nolinkurl{paperI_main}\hspace{0.45em}\nolinkurl{(G}\hspace{0.45em}\nolinkurl{:}\hspace{0.45em}\nolinkurl{Split)}\hspace{0.45em}\nolinkurl{:}\par
\noindent \hspace{1.8em}\nolinkurl{G.Phi}\hspace{0.45em}\nolinkurl{≤}\hspace{0.45em}\nolinkurl{(G.n}\hspace{0.45em}\nolinkurl{:}\hspace{0.45em}\nolinkurl{ℝ)}\hspace{0.45em}\nolinkurl{^}\hspace{0.45em}\nolinkurl{2}\hspace{0.45em}\nolinkurl{/}\hspace{0.45em}\nolinkurl{6}\hspace{0.45em}\nolinkurl{+}\hspace{0.45em}\nolinkurl{(G.n}\hspace{0.45em}\nolinkurl{:}\hspace{0.45em}\nolinkurl{ℝ)}\par
}
\end{quote}

Aquí \nolinkurl{Split} empaqueta los datos de un grafo split usados en la
demostración: una clique de tamaño \nolinkurl{p}, un tipo de vértices
independientes y los vecindarios de esos vértices dentro de la clique.
El campo \nolinkurl{n} es el número total de vértices. La cantidad
\nolinkurl{Phi} se define como
\texttt{\textbar{}E\textbar{}\ −\ 2·nu3star}, donde \nolinkurl{nu3star} es
el supremo de los valores factibles de empaquetamiento fraccional de
triángulos. En otros artículos de la serie este déficit fraccional se
escribe \(\Phi^*\), mientras que \(\Phi\), sin estrella, se reserva para
el déficit integral \(|E(G)|-2\nu_3(G)\); el identificador Lean se
conserva para coincidir con la declaración registrada. Con estas
definiciones, \nolinkurl{paperI_main_sharp} es exactamente

\[
|E(G)|-2\nu_3^*(G)\le \frac{n^2}{6}+\frac{n}{2},
\]

la desigualdad enunciada en el Teorema 1.1 (declaración
\nolinkurl{paperI_main_sharp}). La declaración compañera
\nolinkurl{paperI_main} registra el término aditivo más débil \(n\).

El informe de build suministrado registra cero \nolinkurl{sorry}. El
informe de axiomas suministrado es:

\begin{quote}
{\LeanSymbols\fontsize{8}{10}\selectfont
\noindent \nolinkurl{PaperI.Split.paperI_main_sharp}\hspace{0.45em}\nolinkurl{depends}\hspace{0.45em}\nolinkurl{on}\hspace{0.45em}\nolinkurl{axioms:}\hspace{0.45em}\nolinkurl{[propext,}\hspace{0.45em}\nolinkurl{Classical.choice,}\hspace{0.45em}\nolinkurl{Quot.sound]}\par
\noindent \nolinkurl{PaperI.paperI_main}\hspace{4.05em}\nolinkurl{depends}\hspace{0.45em}\nolinkurl{on}\hspace{0.45em}\nolinkurl{axioms:}\hspace{0.45em}\nolinkurl{[propext,}\hspace{0.45em}\nolinkurl{Classical.choice,}\hspace{0.45em}\nolinkurl{Quot.sound]}\par
\noindent \nolinkurl{PaperI.assembly_sharp}\hspace{2.7em}\nolinkurl{depends}\hspace{0.45em}\nolinkurl{on}\hspace{0.45em}\nolinkurl{axioms:}\hspace{0.45em}\nolinkurl{[propext,}\hspace{0.45em}\nolinkurl{Classical.choice,}\hspace{0.45em}\nolinkurl{Quot.sound]}\par
\noindent \nolinkurl{PaperI.Split.residual_duality}\hspace{0.45em}\nolinkurl{depends}\hspace{0.45em}\nolinkurl{on}\hspace{0.45em}\nolinkurl{axioms:}\hspace{0.45em}\nolinkurl{[propext,}\hspace{0.45em}\nolinkurl{Classical.choice,}\hspace{0.45em}\nolinkurl{Quot.sound]}\par
}
\end{quote}

Estos son los tres axiomas fundacionales estándar de Lean y Mathlib. No
se usa ningún axioma específico del proyecto.

\textbf{Tabla 2. Perímetro de la formalización congelada (freeze Lean
v1.2).} El snapshot congelado de Lean 4 / Mathlib v4.28.0 registra el
núcleo demostrado usando únicamente los axiomas fundacionales estándar
de Lean y ningún axioma matemático propio del proyecto. Una compilación
externa en clean room y una comprobación de axiomas por teorema
reprodujeron el archivo exacto.

{\def\LTcaptype{none} % do not increment counter
\begin{longtable}[]{@{}
  >{\raggedright\arraybackslash}p{(\linewidth - 4\tabcolsep) * \real{0.3333}}
  >{\raggedright\arraybackslash}p{(\linewidth - 4\tabcolsep) * \real{0.3333}}
  >{\raggedright\arraybackslash}p{(\linewidth - 4\tabcolsep) * \real{0.3333}}@{}}
\toprule\noalign{}
\begin{minipage}[b]{\linewidth}\raggedright
Nodo
\end{minipage} & \begin{minipage}[b]{\linewidth}\raggedright
Enunciado
\end{minipage} & \begin{minipage}[b]{\linewidth}\raggedright
Huella axiomática
\end{minipage} \\
\midrule\noalign{}
\endhead
\bottomrule\noalign{}
\endlastfoot
\nolinkurl{paperI_main_sharp} (Teorema 1.1) &
\(\lvert E(G)\rvert-2\nu_3^*(G)\le n^2/6+n/2\) & \nolinkurl{propext},
\nolinkurl{Classical.choice}, \nolinkurl{Quot.sound} \\
\nolinkurl{paperI_main} (cota compañera más débil) &
\(\lvert E(G)\rvert-2\nu_3^*(G)\le n^2/6+n\) & \nolinkurl{propext},
\nolinkurl{Classical.choice}, \nolinkurl{Quot.sound} \\
\nolinkurl{lpVal_concave}, \nolinkurl{lpVal_ge_min} (interfaz de §5) &
concavidad abstracta / cota inferior por perfiles puros &
\nolinkurl{propext}, \nolinkurl{Classical.choice}, \nolinkurl{Quot.sound} \\
\nolinkurl{lpVal_lipschitz} (interfaz de §5) & estabilidad paramétrica del
valor fraccional & \nolinkurl{propext}, \nolinkurl{Classical.choice},
\nolinkurl{Quot.sound} \\
\nolinkurl{residual_duality} (Prop. B.1) & dualidad finita
empaquetamiento--cobertura & \nolinkurl{propext},
\nolinkurl{Classical.choice}, \nolinkurl{Quot.sound} \\
\end{longtable}
}

\subsubsection{Proposición B.1}\label{proposiciuxf3n-b.1}

\textbf{Nota de formalización.} El freeze v1.2 registra la Proposición
B.1 como la declaración Lean \nolinkurl{PaperI.Split.residual_duality},
una instancia del teorema general de dualidad fuerte finita
\nolinkurl{covering_packing_duality}. El artefacto registra este último
como probado a partir de un lema de Farkas finito, usando el carácter
cerrado de un cono convexo finitamente generado y el teorema de
separación por hiperplanos para conos convexos cerrados. Por tanto, la
dualidad usada aquí se deriva dentro del desarrollo formal y no se
postula como un axioma específico del proyecto. Schrijver sigue siendo
la referencia clásica para el argumento escrito. Las fuentes
correspondientes, el log de build y el informe de axiomas se entregan en
\nolinkurl{PAPER_I/05_formalization/lean_v1.2_freeze}.

\subsubsection{1. Introducción}\label{introducciuxf3n}

\textbf{Notación de la serie.} En las comparaciones entre artículos, el
déficit fraccional de Paper I se escribe
\(\Phi^*(G)=|E(G)|-2\nu_3^*(G)\). Paper II denota la misma cantidad por
\(\Phi_\tau(G)=|E(G)|-2\tau_3^*(G)\), mientras que Paper III reserva
\(\Phi(G)\), sin estrella, para el déficit integral
\(|E(G)|-2\nu_3(G)\). El identificador Lean \nolinkurl{Phi} del artefacto
de Paper I se conserva por trazabilidad.

\subsubsection{Observación (concavidad sobre un politopo
fijo)}\label{observaciuxf3n-concavidad-sobre-un-politopo-fijo}

La reducción que acabamos de usar es un caso de un hecho puramente
convexo-analítico, que formulamos aquí de manera abstracta. Fijemos un
politopo factible que no dependa del perfil y una función objetivo que
dependa afínmente de la masa del perfil \(\mu=(\mu_S)\). Como el valor
óptimo de un programa lineal sobre un politopo fijo es cóncavo respecto
de su vector objetivo, el valor en un perfil mixto domina la combinación
convexa correspondiente de los valores de los perfiles puros y, por
tanto, también su mínimo: \[
M_b(\kappa_\mu)\ \ge\ \sum_S \tfrac{\mu_S}{q}\,M_b(\kappa^{S,q})\ \ge\ \min_S M_b(\kappa^{S,q}).
\] Esto se registra en el desarrollo formal como \nolinkurl{lpVal_concave}
y \nolinkurl{lpVal_ge_min}. Una estimación complementaria,
\nolinkurl{lpVal_lipschitz}, acota la variación del valor fraccional
cuando se perturban el baseline o la distribución de perfiles. Está
pensada para variantes robustas y casi split. Ambas se formulan para un
politopo factible fijo con una función objetivo que varía afínmente;
ninguna requiere simetría, que solo interviene después, sobre un único
perfil puro, en la Sección 6.

\subsection{Referencias}\label{referencias}

{[}1{]} P. Erdős, E. T. Ordman, and Y. Zalcstein, ``Clique partitions of
chordal graphs,'' \emph{Combinatorics, Probability and Computing}
\textbf{2} (1993), no. 4, 409--415.

{[}2{]} G.-T. Chen, P. Erdős, and E. T. Ordman, ``Clique partitions of
split graphs,'' in \emph{Combinatorics, Graph Theory, Algorithms and
Applications} (Beijing, 1993), World Scientific, 1994, pp.~21--30, copia
alojada por uno de los autores:
\url{https://ordman.net/MathResearch/CEOClique_Parts.pdf}, consultada el
21 de agosto de 2026.

{[}3{]} T. F. Bloom, ``Erdős Problem \#81,'' \emph{Erdős Problems},
\url{https://www.erdosproblems.com/81}, consultado el 21 de agosto de
2026.

{[}4{]} A. Schrijver, \emph{Theory of Linear and Integer Programming},
Wiley-Interscience Series in Discrete Mathematics and Optimization, John
Wiley \& Sons, Chichester, 1986.

{[}5{]} M. S. Cavers, \emph{Clique Partitions and Coverings of Graphs},
M.Math. essay, University of Waterloo, 2005.

{[}6{]} L. de Moura and S. Ullrich, ``The Lean 4 Theorem Prover and
Programming Language,'' in \emph{Automated Deduction -- CADE 28},
Lecture Notes in Computer Science \textbf{12699}, Springer, 2021,
pp.~625--635, \url{https://doi.org/10.1007/978-3-030-79876-5_37}.

{[}7{]} The mathlib Community, ``The Lean Mathematical Library,'' in
\emph{Proceedings of the 9th ACM SIGPLAN International Conference on
Certified Programs and Proofs (CPP 2020)}, ACM, 2020, pp.~367--381,
\url{https://doi.org/10.1145/3372885.3373824}.

\end{document}
